\chapter{Proportionnalité}\label{ch:g7:prop}

L'année 6 a rencontré les tableaux de \omterm{def:g7:prop:table}{proportionnalité}
(\cref{ch:g6:propdata})~; ce chapitre en fait un outil à part
entière~: la règle de trois pour trouver une quatrième valeur, les
pourcentages, et les \omterm{def:g7:prop:scale}{échelles} de cartes et de maquettes. La vitesse,
la plus célèbre de toutes les \omterm{def:g7:prop:table}{proportionnalités}, est étudiée au
\cref{ch:g8:speed}.

\section{Reconnaître et compléter des tableaux}

\begin{definition}[Tableau de proportionnalité]\label{def:g7:prop:table}
Un tableau à deux lignes est un \emph{tableau de
proportionnalité}\index{proportionnalité} lorsque les nombres de la
seconde ligne sont ceux de la première multipliés par un nombre fixe,
le \emph{coefficient}. De façon équivalente~: tous les \omterm{def:g3:division:remainder}{quotients} de
colonne $\frac{\text{bas}}{\text{haut}}$ sont égaux.
\end{definition}

\begin{example}\label{ex:g7:prop:recognize}
\begin{center}
\renewcommand{\arraystretch}{1.3}
\begin{tabular}{l|ccc}
haut & $4$ & $10$ & $14$ \\
\hline
bas & $6$ & $15$ & $21$ \\
\end{tabular}
\end{center}
\omterm{def:g3:division:remainder}{Quotients}~: $\frac64 = 1.5$, $\frac{15}{10} = 1.5$,
$\frac{21}{14} = 1.5$~: proportionnel, coefficient $1.5$.
\end{example}

\begin{theorem}[Règle de trois]\label{thm:g7:prop:cross}
Dans un \omterm{def:g7:prop:table}{tableau de proportionnalité} de colonnes
$\begin{pmatrix} a \\ b \end{pmatrix}$ et
$\begin{pmatrix} c \\ d \end{pmatrix}$, les \omterm{def:g2:mult:def}{produits} en croix sont
égaux~:
\[
a \times d = b \times c .
\]
Par conséquent la quatrième valeur se calcule à partir des trois
autres~: $d = \dfrac{b \times c}{a}$.
\end{theorem}

\begin{proof}
Soit $k$ le coefficient~: $b = ka$ et $d = kc$. Alors
\[
a \times d = a \times kc = k \times ac,
\qquad
b \times c = ka \times c = k \times ac :
\]
les deux \omterm{def:g2:mult:def}{produits} en croix valent $k \times ac$. Diviser $ad = bc$
par $a$ donne la formule pour $d$.
\end{proof}

\begin{example}[Quatrième proportionnelle]\label{ex:g7:prop:fourth}
Si $7$ livres identiques pèsent $2.8$ kg, combien pèsent $12$ tels
livres~? Tableau et règle de trois~:
\begin{center}
\renewcommand{\arraystretch}{1.3}
\begin{tabular}{l|cc}
livres & $7$ & $12$ \\
\hline
kg & $2.8$ & $m$ \\
\end{tabular}
\end{center}
\[
7 \times m = 2.8 \times 12
\quad\Longrightarrow\quad
m = \frac{2.8 \times 12}{7} = \frac{33.6}{7} = 4.8 \text{ kg}.
\]
Vérification par l'unité~: un livre pèse $0.4$ kg, douze pèsent
$4.8$ kg.
\end{example}

\section{Pourcentages}

\begin{method}[Appliquer et trouver des pourcentages]\label{met:g7:prop:percent}
\begin{enumerate}
  \item \emph{Appliquer} $t\,\%$~: multiplier par $\frac{t}{100}$
        ($12\,\%$ de $350$ valent $0.12 \times 350 = 42$)~;
  \item \emph{trouver} le pourcentage qu'une partie représente~:
        calculer $\frac{\text{partie}}{\text{tout}}$ et réécrire sur
        $100$ ($21$ sur $60$~: $\frac{21}{60} = \frac{35}{100} =
        35\,\%$)~;
  \item toujours se demander~: \emph{un pourcentage de quoi~?} Le
        tout de référence compte.
\end{enumerate}
\end{method}

\begin{example}\label{ex:g7:prop:percent}
Dans une école, $180$ des $400$ élèves sont des filles.
Pourcentage~:
\[
\frac{180}{400} = \frac{180 \div 4}{400 \div 4} = \frac{45}{100}
= 45\,\% .
\]
Dans l'autre sens~: combien sont les $55\,\%$ de garçons~?
$0.55 \times 400 = 220$. Vérification~: $180 + 220 = 400$.
\end{example}

\section{Échelles}

\begin{definition}[Échelle]\label{def:g7:prop:scale}
Sur une carte ou une maquette à \emph{échelle}\index{échelle}
$\frac{1}{n}$ (écrite $1 : n$), les distances sur la carte sont
\omterm{def:g6:propdata:def}{proportionnelles} aux distances réelles, avec le coefficient
$\frac1n$~:
\[
\text{distance carte} = \frac{1}{n} \times \text{distance réelle},
\]
toutes deux mesurées dans la \emph{même unité}.
\end{definition}

\begin{example}\label{ex:g7:prop:scale}
Sur une carte au $1 : 50\,000$, deux villes sont distantes de $7$ cm.
Distance réelle~:
\[
7 \times 50\,000 = 350\,000 \text{ cm} = 3\,500 \text{ m} = 3.5
\text{ km}.
\]
Inversement, un sentier de $12$ km mesure sur la carte
$12$ km $= 1\,200\,000$ cm, donc $1\,200\,000 \div 50\,000 = 24$ cm.
\end{example}

\begin{omfigure}
\begin{tikzpicture}
\begin{axis}[omaxis,
  width=8.4cm, height=5.6cm,
  xmin=0, xmax=8.6, ymin=0, ymax=34,
  xlabel={carte (cm)}, ylabel={réel (km)},
  xtick={2,4,6,8}, ytick={10,20,30}]
  \addplot[omDef, thick, domain=0:8.2] {4*x};
  \addplot[omDef, only marks, mark=*, mark size=1.5pt] coordinates
    {(2,8) (4,16) (6,24) (8,32)};
  \draw[dashed, omThm] (axis cs:7,0) -- (axis cs:7,28)
    -- (axis cs:0,28);
\end{axis}
\end{tikzpicture}

{\small Une \omterm{def:g7:prop:scale}{échelle} est une \omterm{def:g7:prop:table}{proportionnalité}~: distances carte contre
distances réelles s'alignent sur une droite passant par l'origine (ici
$1$ cm $\leftrightarrow$ $4$ km~; les tirets lisent $7$ cm
$\leftrightarrow$ $28$ km).}
\end{omfigure}

\begin{remark}
Les graphiques sont le test le plus rapide de \omterm{def:g7:prop:table}{proportionnalité}
(\cref{ch:g6:propdata})~: des points sur une droite \emph{passant par
l'origine} signifient des quantités \omterm{def:g6:propdata:def}{proportionnelles}~; une droite
manquant l'origine (comme le prix d'un taxi avec prise en charge)
signifie \emph{non} proportionnel, même si c'est une droite. Les
fonctions affines, au \cref{ch:g9:linfunc}, rendront cela précis.
\end{remark}

\section{Exercices}

\begin{exercise}[$\star$]\label{exo:g7:prop:1}
Quels tableaux sont des tableaux de \omterm{def:g7:prop:table}{proportionnalité}~? Donner le
coefficient lorsqu'il existe.
\begin{center}
\renewcommand{\arraystretch}{1.2}
\begin{tabular}{l|ccc}
 & $3$ & $7$ & $11$ \\
\hline
 & $7.5$ & $17.5$ & $27.5$ \\
\end{tabular}
\qquad
\begin{tabular}{l|ccc}
 & $2$ & $5$ & $9$ \\
\hline
 & $6$ & $15$ & $28$ \\
\end{tabular}
\end{center}
\end{exercise}

\begin{exercise}[$\star$]\label{exo:g7:prop:2}
Compléter par la règle de trois, en écrivant le calcul~: $5$ kg de
pommes de terre coûtent $6$ euros~; combien coûtent $8$ kg~?
\end{exercise}

\begin{exercise}[$\star$]\label{exo:g7:prop:3}
Une imprimante imprime $36$ pages en $3$ minutes. Au même rythme,
combien de pages en $5$ minutes~? Combien de temps pour $96$ pages~?
\end{exercise}

\begin{exercise}[$\star$]\label{exo:g7:prop:4}
Calculer~: $30\,\%$ de $250$~; \quad $15\,\%$ de $60$~; \quad
$t\,\%$ tel que $t\,\%$ de $80$ vaut $20$.
\end{exercise}

\begin{exercise}[$\star$]\label{exo:g7:prop:5}
Sur $25$ tirs, une basketteuse a marqué $18$. Quel est son
pourcentage de réussite~? (Réécrire la \omterm{def:g6:fractions:def}{fraction} sur $100$.)
\end{exercise}

\begin{exercise}[$\star$]\label{exo:g7:prop:6}
Sur une carte au $1 : 25\,000$, un chemin mesure $9$ cm. Quelle est
sa longueur réelle en km~? Un lac mesure $2$ km de long~: quelle
longueur sur la carte~?
\end{exercise}

\begin{exercise}[$\star$]\label{exo:g7:prop:7}
Une voiture miniature est construite à l'\omterm{def:g7:prop:scale}{échelle} $1 : 43$. La
voiture réelle mesure $4.3$ m de long. Quelle longueur a la
miniature, en cm~?
\end{exercise}

\begin{exercise}[$\star\star$]\label{exo:g7:prop:8}
Une recette utilise $240$ g de chocolat pour $8$ parts. Aline dispose
de $300$ g de chocolat. Pour combien de parts cela suffit-il (nombre
entier de parts)~?
\end{exercise}

\begin{exercise}[$\star\star$]\label{exo:g7:prop:9}
Le prix d'une veste passe de $80$ à $60$ euros.
\begin{enumerate}
  \item Quelle est la réduction en euros~? En pour cent du prix
        initial~?
  \item Plus tard le prix remonte de $60$ à $80$ euros. Expliquer
        pourquoi cette hausse n'est \emph{pas} de $25\,\%$ ---
        calculer le pourcentage d'augmentation correct.
\end{enumerate}
\end{exercise}

\begin{exercise}[$\star\star$]\label{exo:g7:prop:10}
Deux lignes d'un tableau sont \omterm{def:g6:propdata:def}{proportionnelles}~:
\begin{center}
\renewcommand{\arraystretch}{1.2}
\begin{tabular}{l|ccc}
 & $4$ & $x$ & $10$ \\
\hline
 & $y$ & $10.5$ & $17.5$ \\
\end{tabular}
\end{center}
Trouver le coefficient, puis $x$ et $y$.
\end{exercise}

\begin{exercise}[$\star\star\star$]\label{exo:g7:prop:11}
Un photocopieur réduit les documents à $80\,\%$ de leur taille. Un
\omterm{def:g6:lines:objects}{segment} de $10$ cm est copié, puis la \emph{copie} est copiée à
nouveau avec le même réglage.
\begin{enumerate}
  \item Quelle longueur a le \omterm{def:g6:lines:objects}{segment} après une copie~? Après deux~?
  \item Pourquoi la réponse après deux copies n'est-elle pas $60\,\%$
        de l'original~? Quel pourcentage unique correspond à deux
        copies~?
\end{enumerate}
\end{exercise}

\section{Problème~: mesurer la Terre avec un bâton}

\begin{problem}[{Devoir maison --- les ombres forment une
proportionnalité, et comment Ératosthène a calculé la circonférence
de la Terre en 240 av. J.-C.}]
\label{pb:g7:prop:1}
Il y a vingt-deux siècles, sans télescope, sans satellite et sans
calculatrice, le bibliothécaire d'Alexandrie a mesuré la Terre ---
avec un bâton, un puits, une caravane de chameaux et les
mathématiques de ce chapitre. Ce problème refait son chemin~: la
\omterm{def:g7:prop:table}{proportionnalité} des ombres d'abord, puis le calcul célèbre
lui-même, et enfin ce que devient sa réponse à l'\omterm{def:g7:prop:scale}{échelle} humaine.

\medskip
\noindent\textbf{Partie I --- L'ombre d'un bâton.}
Le Soleil est si loin que ses \omterm{def:g3:shapes:shapes}{rayons} nous parviennent \omterm{def:g6:lines:perp}{parallèles}
entre eux. À un instant donné, tous les objets verticaux et leurs
ombres sont donc proportionnels.
\begin{enumerate}
  \item À midi, un bâton vertical de $1$ m projette une ombre de
        $0.4$ m, et l'ombre d'un arbre mesure $3.2$ m. Quelle est la
        hauteur de l'arbre (\cref{thm:g7:prop:cross})~?
  \item Au même instant, une tour projette une ombre de $26$ m~:
        quelle est sa hauteur~? Et quelle est la longueur de l'ombre
        d'un enfant mesurant $1.5$ m~?
  \item Expliquer en une phrase pourquoi tous ces calculs ne valent
        qu'au \emph{même instant} de la journée.
  \item La légende raconte que Thalès a mesuré la grande \omterm{def:g5:solids:def}{pyramide} en
        attendant l'instant où \emph{sa propre ombre avait exactement
        la longueur de sa taille}. Que vaut le coefficient de
        \omterm{def:g7:prop:table}{proportionnalité} à cet instant, et quelle est la hauteur de
        la \omterm{def:g5:solids:def}{pyramide} si l'extrémité de son ombre se trouve à $147$ m
        du point du sol situé à la verticale de son \omterm{def:g5:solids:def}{sommet}~?
  \item Résumer la partie I~: à un instant donné, quelle grandeur
        est la même pour le bâton, l'arbre, la tour et la \omterm{def:g5:solids:def}{pyramide}
        (\cref{def:g7:prop:table})~?
\end{enumerate}

\medskip
\noindent\textbf{Partie II --- Le calcul d'Ératosthène.}
Ératosthène connaissait deux faits. Dans la ville de Syène, à midi,
le jour du solstice d'été, le Soleil se tenait \emph{exactement à la
verticale}~: la lumière atteignait le fond des puits les plus
profonds, et les bâtons verticaux ne projetaient aucune ombre. À
Alexandrie, $800$ km plein nord, au même instant, un bâton vertical
projetait \emph{bel et bien} une ombre, et les \omterm{def:g3:shapes:shapes}{rayons} du Soleil
faisaient un angle de $7.2^\circ$ avec la verticale.
\begin{enumerate}[resume]
  \item Expliquer pourquoi ces deux observations réunies prouvent que
        le sol de Syène et le sol d'Alexandrie ne sont pas \omterm{def:g6:lines:perp}{parallèles}
        --- autrement dit, que la surface de la Terre est courbe.
        (Que feraient des \omterm{def:g3:shapes:shapes}{rayons} \omterm{def:g6:lines:perp}{parallèles} à deux bâtons sur une
        Terre plate~?)
  \item Sur un dessin de la Terre ronde éclairée par des \omterm{def:g3:shapes:shapes}{rayons}
        \omterm{def:g6:lines:perp}{parallèles}, les $7.2^\circ$ d'Alexandrie réapparaissent au
        centre de la Terre, comme angle entre les directions des deux
        villes. Faire le dessin et s'en convaincre (la justification
        propre utilise les couples d'angles égaux du
        \cref{ch:g7:triangles}).
  \item Quelle \omterm{def:g6:fractions:def}{fraction} d'un tour complet représentent $7.2^\circ$~?
  \item L'arc de Syène à Alexandrie ($800$ km) correspond à
        $7.2^\circ$~; la circonférence entière correspond à
        $360^\circ$. Poser la \omterm{def:g7:prop:table}{proportionnalité} et calculer la
        circonférence de la Terre.
  \item En déduire le \omterm{def:g4:geometry:circle}{diamètre} de la Terre, avec
        $\pi \approx 3.14$ (\cref{prop:g6:measure:circle}) et en
        arrondissant à la centaine de kilomètres. (La valeur moderne
        est $12\,742$ km --- de combien s'en approchait un homme
        muni d'un bâton en 240 av. J.-C.~?)
\end{enumerate}

\medskip
\noindent\textbf{Partie III --- La Terre à l'\omterm{def:g7:prop:scale}{échelle} humaine.}
\begin{enumerate}[resume]
  \item Un globe est construit à l'\omterm{def:g7:prop:scale}{échelle} $1 : 40\,000\,000$. À
        l'aide de la circonférence trouvée à la question 9, montrer
        que la circonférence du globe vaut exactement $1$ m. Quel est
        son \omterm{def:g4:geometry:circle}{diamètre}, au centimètre près~?
  \item Le mont Blanc s'élève à environ $4.8$ km. Convertir $4.8$ km
        en centimètres, diviser par $40\,000\,000$, et exprimer la
        hauteur de la montagne sur le globe en millimètres. Que dit
        cette réponse sur la véritable «~rugosité~» de la Terre~?
  \item Un marcheur parcourt $40$ km par jour. À ce rythme, combien
        de jours pour la circonférence entière d'Ératosthène --- et
        cela fait à peu près combien d'années~?
  \item L'atmosphère respirable est concentrée dans les $10$
        premiers kilomètres au-dessus du sol, en gros. Quel
        pourcentage du \omterm{def:g3:shapes:shapes}{rayon} terrestre (environ $6\,370$ km) cela
        représente-t-il (\cref{met:g7:prop:percent})~? Arrondir au
        dixième de pour cent.
  \item Les historiens estiment que la valeur annoncée par
        Ératosthène, convertie en unités modernes, valait peut-être
        environ $42\,000$ km. En prenant $40\,000$ km pour valeur
        exacte, calculer son pourcentage d'erreur. Conclure en une
        phrase sur les bâtons, les puits et la \omterm{def:g7:prop:table}{proportionnalité}.
\end{enumerate}
\end{problem}
